\chapter{Risks, Open Items and Next Steps}\label{ch:risks}

\section{Risks}
\begin{itemize}
  \item \textbf{Codec and H750 next to a high-impedance guitar input} on one 56 by 90\,mm board. Mitigation: The Relic's four-layer stack with an unbroken ground plane, the zoning in Chapter~\ref{ch:pcb}, and the buck converter at the DC end away from the input.
  \item \textbf{Chord tracker quality} on distorted or fast playing. Key, Drone and Fifths modes do not depend on it, so the pedal still works if Follow is weak.
  \item \textbf{CPU load} of 24 strings, Jawari and the tracker together has not been measured.
  \item \textbf{Internal flash.} Whether the image fits the H750's 128\,KB is unknown.
  \item \textbf{AK4621EF lifecycle,} already flagged in the Timekeeper's BOM audit.
  \item \textbf{Light pipes} add a hand-assembly step and a part to source.
\end{itemize}

\section{Open Items}
\begin{table}[H]
  \small\renewcommand{\arraystretch}{1.12}
  \caption{Open hardware and product items.}\label{tab:openhw}
  \begin{tabularx}{\textwidth}{@{}>{\raggedright\arraybackslash}p{38mm} L@{}}
    \toprule
    \fmtophead{Item} & \fmtophead{Status} \\
    \midrule
    Toggle row & Keep only the centre toggle, or give up a jack (Section~\ref{sec:toggleconflict}). Then pick the toggle part: a two-position Taiway 100-series toggle on The Alchemist's footprint, or The Alchemist's ON-ON-ON part with a second input pin per toggle. \\
    Schematic confirmations & QUADSPI bank~2 alternate functions, AK4621 unused right input, THS4522 unused channel, OPA2365 noise at 1\,kHz, then ERC. \\
    Wall-hole heights & To be fixed from the NMJ6HCD2 and PJ-063AH drawings and the measured Alchemist and Relic assemblies; the Harp inherits them. \\
    Note display & Whether 12 LEDs are enough, or a small I$^2$C OLED in the same band is worth the window cut. Judged after the firmware prototype, using the Timekeeper's TFT as a stand-in. \\
    Light-pipe fit & Flange diameter against the 5\,mm pitch and the cavity wall. \\
    Expression input & TRS wiring and buffer as the Timekeeper (MCP6001); whether an empty jack can be detected depends on that buffer. \\
    Mono use & The firmware cannot tell that OUT~R is unplugged. Either choose a default spread a mono player can live with, or add a jack-detect line while pins are free. \\
    Power-up state & Effect on, or the last state. \\
    Secondary-layer gesture & Both footswitches together, or holding Bypass. \\
    Code sharing with the Timekeeper & Copy the reused modules with the source commit recorded, or share them as a library. \\
    \bottomrule
  \end{tabularx}
\end{table}

\section{Discrepancies Between the Working Documents}
Writing this report found three small disagreements, to be corrected at the source:
\begin{itemize}
  \item The design brief's latency figure (one 256-sample block, about 6\,ms) does not match the Timekeeper's 64-frame block (about 1.5\,ms); the firmware requirements use the smaller figure.
  \item The firmware requirements put the expression input on PA6 [EXP-1]; the schematic and pin map use PB1, as the Timekeeper does.
  \item The firmware requirements name the footswitches J105 and J106; the schematic designators are \req{J405} (Bypass) and \req{J407} (Hold).
\end{itemize}

\section{Order of Work}
\begin{enumerate}
  \item \textbf{Firmware first on the Timekeeper board.} String bank with a fixed tuning, then Jawari, then the tuning manager in Key, Drone and Fifths modes. Measure the CPU load and settle the TBC ranges by ear.
  \item Chord tracker and Follow mode, offline against reference recordings, then on the board.
  \item Open the schematic in KiCad and run ERC. Decide the toggle conflict and close the confirmation items.
  \item \emph{Update PCB from Schematic}. Place the LED strip and check the light-pipe fit; place the MCU and flash in the centre band, the codec under the pots, the buffers at the jack ends and the buck on the tab. \textbf{Go/no-go for Plan~A.}
  \item GND vias per The Relic's rule, routing (F.Cu short, B.Cu long, In2 power), DRC with The Relic's rules.
  \item Adapt The Relic's BOM exporter, cost at 100, and produce the Tayda drill schedule from the face table.
  \item Fix the wall-hole heights from the measured Alchemist and Relic assembly.
  \item Port the firmware to \texttt{BOARD\_HARP}, then bench measurements and the production test mode.
\end{enumerate}

\section{Source Documents}
\begin{itemize}
  \item \texttt{docs/design-brief.md}: what the pedal is, controls, bypass, cost and decisions.
  \item \texttt{docs/pcb-plan.md}: board plan, face schedule, zoning, schematic sheets, pin map and input-stage checks.
  \item \texttt{docs/firmware-requirements.md}: numbered firmware requirements, Draft~A.
  \item \texttt{kicad/the\_harp/} and \texttt{kicad/README.md}: the KiCad project.
  \item \texttt{scripts/harp\_schematic/}: the schematic generator and net checker.
  \item 25-Z01-0001 DSP Development Board: the Timekeeper schematic, firmware, input noise review and input-chain simulation.
  \item 26-A03-0003 The Relic and 26-A02-0001 The Alchemist: board outline, face standard, jacks, relay bypass and layout rules.
\end{itemize}
